\chapter{超空间与超对称动力学的构造}

\lectureoverview{我们从超对称能够作用的最小系统出发重建超对称：一个玻色振子和一个能量相同的费米振子。它们的二阶方程与一阶方程之间的对比引出了超荷及其代数。随后我们将该代数重新解释为具有普通坐标 (t) 和格拉斯曼坐标 \(\theta,\bar\theta\) 的空间上的几何。超场将玻色子、费米子和一个辅助场打包成一个单一对象；别列津积分提取出不变量的普通拉格朗日量；协变导数使与对称性相容的约束成为可能；手征约束则提供了第一个重要的例子。本讲刻意暴露这套机制在符号和排序上的困难，最后说明为什么对称性本身比一份费曼规则目录更为基本。}

\section{回到单模系统}

让我们再次从静止的费米子和玻色子开始。我们抑制掉所有空间坐标和所有动量标签。在给定模式中，一个量子要么存在要么不存在，而产生和湮灭算符就足以描述它。记 \(a^\dagger,a\) 为玻色模式，\(c^\dagger,c\) 为费米模式。它们满足
\begin{align}
[a,a^\dagger]&=1,
&
\{c,c^\dagger\}&=1,
&
c^2=(c^\dagger)^2&=0.
\end{align}
如果这两个量子具有相同的能量 \(m\)，那么，除去一个无关的真空常数，
\begin{equation}
H=m\left(a^\dagger a+c^\dagger c\right).
\label{eq:l7-one-mode-hamiltonian}
\end{equation}
数算符 \(a^\dagger a\) 可以取任意非负整数值，而 \(c^\dagger c\) 只能为零或一。这就是这个玩具模型的全部运动学。

这种简化是策略性的。我们想要理解那种在不改变能量的情况下移除一个玻色子并插入一个费米子、或者反过来进行的操作。然而，在构造该操作之前，我们先把同一个系统从算符语言翻译成拉格朗日量语言。

\section{玻色振子}

对于玻色子，引入一个实坐标 \(\phi(t)\)。在没有空间变化的情况下，自由标量场变成了一个普通的谐振子。黑板上写下的第一个表达式是
\[
\frac{1}{2}\dot\phi^{,2}+\frac{1}{2}m^2\phi^2.
\]
那个表达式是动能加势能，而不是拉格朗日量。讲演中立即修正了符号：
\begin{equation}
L_B=\frac{1}{2}\dot\phi^{,2}-\frac{1}{2}m^2\phi^2.
\label{eq:l7-boson-lagrangian}
\end{equation}

\begin{figure}[H]
\centering
\includegraphics[width=0.62\textwidth]{lecture_07_figure_02.png}
\lectureframe{00:03:25.431}
\caption{玻色拉格朗日量之上的单模哈密顿量。该画面保留了口头修正为 \(T-V\) 之前的初始正号。}
\end{figure}

\begin{classroomqa}
\audiencequestion{所展示的振子拉格朗日量中的符号有什么问题？\lecturetimestamp{00:04:00}}
\lecturerresponse{能量是 \(T+V\)，但拉格朗日量是 \(T-V\)。因此势能项必须带一个负号。这一修正将方程从指数行为变为所期望的谐振子振荡。}
\end{classroomqa}

的确，欧拉--拉格朗日方程给出
\begin{align}
\frac{d}{dt}\frac{\partial L_B}{\partial\dot\phi}
&=\frac{\partial L_B}{\partial\phi},\\
\ddot\phi&=-m^2\phi,
\end{align}
或
\begin{equation}
\ddot\phi+m^2\phi=0.
\label{eq:l7-boson-eom}
\end{equation}
其频率是 \(m\)，而不是 \(m^2\)。对这个振子进行量子化会产生由式~\eqref{eq:l7-one-mode-hamiltonian}中的 \(a^\dagger\) 和 \(a\) 所描述的同一阶梯态。

\section{为什么费米子作用量是一阶的}

费米拉格朗日量看起来更简单，因为它只含有一个时间导数。原因在一个刻意朴素的尝试中已经可见。对于一个普通的对易变量，
\begin{equation}
L_{\mathrm{trial}}=\phi\dot\phi
=\frac{1}{2}\frac{d}{dt}\phi^2.
\end{equation}
它是一个全导数。其欧拉--拉格朗日方程只给出 \(\dot\phi=\dot\phi\)，所以每一条可微路径都是允许的。从作用量中移除一个全导数会改变端点数据，但不改变体运动方程。

现在把变量改名为 \(\psi\)，并宣布它是格拉斯曼奇的。\(\psi\) 和 \(\dot\psi\) 都反交换。在固定选择左变分的情况下，重新排列奇因子会引入一个负号，因此同样的形式尝试给出
\begin{equation}
\frac{d}{dt}\frac{\partial L_{\mathrm{trial}}}{\partial\dot\psi}
=\dot\psi,
\qquad
\frac{\partial L_{\mathrm{trial}}}{\partial\psi}
=-\dot\psi.
\end{equation}
因此方程为
\begin{equation}
\dot\psi=-\dot\psi
\quad\Longrightarrow\quad
\dot\psi=0.
\label{eq:l7-constant-fermion}
\end{equation}
对于对易坐标曾是同义反复的表达式，对于反交换坐标却获得了动力学内容。对于振子来说，它的内容太多了，因为它只允许一个常数，但这正是我们所需要的线索。

\begin{classroomqa}
\audiencequestion{解难道就只是 \(\psi\) 为常数吗？\lecturetimestamp{00:13:57}}
\lecturerresponse{是的。式~\eqref{eq:l7-constant-fermion}是一个真正的一阶微分方程，其解是一个任意的格拉斯曼值常数。一个质量项将把这个静态解变成一个振荡解。}
\end{classroomqa}

显而易见的尝试 \(m\psi^2\) 之所以失败，是因为 \(\psi^2=0\)。我们需要两个实奇变量 \(\psi_1,\psi_2\)，或者等价地，一个复变量及其共轭。讲座在黑板上完整推演了双实变量形式，反复核对符号和因子二。一个按惯例归一化的复形式为
\begin{equation}
L_F=\frac{i}{2}
\left(\bar\psi\dot\psi-\dot{\bar\psi}\psi\right)
-m\bar\psi\psi.
\label{eq:l7-fermion-lagrangian}
\end{equation}
在相差一个全导数意义下，这就是 \(i\bar\psi\dot\psi-m\bar\psi\psi\)。对 \(\bar\psi\) 变分得到
\begin{equation}
i\dot\psi=m\psi,
\qquad
\psi(t)=\psi(0)e^{-imt}.
\label{eq:l7-fermion-eom}
\end{equation}
共轭方程具有相反的相位。因此费米子模式以能量 \(m\) 振荡，正如玻色子模式一样，尽管其方程是一阶而非二阶的。

\editorialnote{方程 ~\eqref{eq:l7-fermion-lagrangian} 为探索性的双实变量计算固定了一种标准排序和归一化。讲座本身插入了因子 \(i\)，调整了一个因子二，并留下一个实分量符号未定。这些改动并不改变其核心结论：费米子动力学是一阶的，质量线性地出现，而一个复费米子提供非零的双线性量。}

这是一个熟悉相对论性对比的零维缩影。Klein--Gordon 方程对时间是二阶的，而 Dirac 方程是一阶的。在空间被压掉后，玻色子是一个谐振子，费米子是一个复 Grassmann 振子。

\begin{classroomqa}
\audiencequestion{\(\psi\)、\(\dot\psi\) 及其共轭的书写顺序有关系吗？\lecturetimestamp{00:18:58}}
\lecturerresponse{中间公式确实依赖于顺序，因为奇因子反交换。物理理论并不依赖于任意的书写偏好，只要在整个过程中对左变分或右变分、共轭和积分使用单一约定。交换一对奇因子会改变一个符号，而这个符号必须在其他所有地方一致地改变。}
\end{classroomqa}

结合式 ~\eqref{eq:l7-boson-lagrangian} 和 \eqref{eq:l7-fermion-lagrangian}，给出了式 ~\eqref{eq:l7-one-mode-hamiltonian} 所表示的同一配对的一个拉格朗日量描述。我们现在可以来识别它们之间的对称性了。

\section{超荷及其代数}

一个移除费米子并插入玻色子的算符是
\begin{equation}
Q^\dagger=\sqrt{2m}\,a^\dagger c,
\end{equation}
而其共轭移除玻色子并插入费米子，
\begin{equation}
Q=\sqrt{2m}\,a c^\dagger.
\end{equation}
归一化的选取使得代数取标准形式。利用 \(cc^\dagger=1-c^\dagger c\)、\(aa^\dagger=a^\dagger a+1\)，以及玻色子算符与费米子算符彼此对易这一事实，
\begin{align}
\{Q,Q^\dagger\}
&=2m\left(ac^\dagger a^\dagger c+a^\dagger c a c^\dagger\right)\notag\\
&=2m\left[(a^\dagger a+1)c^\dagger c
+a^\dagger a(1-c^\dagger c)\right]\notag\\
&=2m\left(a^\dagger a+c^\dagger c\right)=2H.
\label{eq:l7-superalgebra}
\end{align}
玻色子算符可以穿过费米子算符。限制费米子占据数的同一幂零性给出
\begin{equation}
Q^2=(Q^\dagger)^2=0.
\end{equation}
最后，
\begin{equation}
[H,Q]=[H,Q^\dagger]=0.
\label{eq:l7-conserved-supercharges}
\end{equation}
每个超荷以相反方向改变玻色子数和费米子数，因此当频率相等时，两个能量变化相消。

\begin{classroomqa}
\audiencequestion{\([H,Q]\) 是什么？答案正比于 \(Q\) 吗？\lecturetimestamp{00:25:10}}
\lecturerresponse{不。答案为零。单独的创生或湮灭算符在与 \(H\) 对易时会获得其能量，但 \(Q\) 中的玻色子贡献和费米子贡献相消。这一相消正是等能量条件。}
\end{classroomqa}

\begin{classroomqa}
\audiencequestion{那个对易子不是已经写在黑板上了吗？\lecturetimestamp{00:25:45}}
\lecturerresponse{是的。重新发现它是有用的，因为这迫使我们解释这一关系，而不仅仅是照抄：与哈密顿量对易意味着该变换是一种对称性，且其生成元守恒。}
\end{classroomqa}

形容词\emph{荷}反映了这一守恒定律。然而，我们不应把 \(Q\) 当作附着于态上的普通可同时测量的数。它是 Grassmann 奇的、幂零的，并把一个玻色子态映射为一个费米子态。有用的陈述是：它是在简并超多重态内作用的一个守恒生成元。
\editorialnote{这一措辞细化了课堂上用能量或电荷标记态的类比。单个非厄米奇超荷不是具有数值量子标签的普通可观测量。它们的守恒作用以及整个超代数的表示提供了正确的分类。}

\section{作为导数的生成元}

这个代数很紧凑，但仅凭它本身并不能告诉我们如何构造相互作用理论。下一步是几何的。熟悉的对称性生成元已经作为微分算符作用。按照讲座的约定，
\begin{equation}
H=i\frac{\partial}{\partial t},
\qquad
P=-i\frac{\partial}{\partial x}.
\end{equation}

\begin{figure}[H]
\centering
\includegraphics[width=0.50\textwidth]{lecture_07_figure_03.png}
\lectureframe{00:28:12.720}
\caption{哈密顿量被表示为时间平移的生成元。}
\end{figure}

角动量同样是一个坐标乘以一个导数。如果超荷是奇生成元，那么自然要问它们能否用关于奇坐标的导数来表示。因此我们把时间轴扩大为超空间
\begin{equation}
(t,\theta,\bar\theta),
\qquad
\theta^2=\bar\theta^2=0,
\qquad
\theta\bar\theta=-\bar\theta\theta.
\end{equation}
这个空间上的场依赖于一个普通时间和一对共轭的格拉斯曼坐标。

本讲刻意揭示符号有多么容易出错。来自复共轭的符号、来自移动一个奇参数的符号，以及来自选择左导数而非右导数的符号，是三种不同的操作。与其保留彼此不相容的黑板阶段，我们现在固定一种约定：所有奇导数都是左导数，并且
\begin{equation}
q=\partial_\theta+i\bar\theta\partial_t,
\qquad
\bar q=\partial_{\bar\theta}+i\theta\partial_t.
\label{eq:l7-differential-supercharges}
\end{equation}
于是
\begin{equation}
q^2=\bar q^{,2}=0,
\qquad
\{q,\bar q\}=2i\partial_t=2H.
\label{eq:l7-differential-algebra}
\end{equation}
因子2来自两个交叉项：\(\{\partial_\theta,\theta\}=1\) 和 \(\{\partial_{\bar\theta},\bar\theta\}=1\)。如果生成元只包含 \(\partial_\theta\) 和 \(\partial_{\bar\theta}\)，它们的反对易子就会为零，无法重现哈密顿量。

\begin{classroomqa}
\audiencequestion{\(q\) 和 \(\bar q\) 中的两个时间导数项应该带相同的符号吗？\lecturetimestamp{00:31:01}}
\lecturerresponse{它们必须被一致地选取，使得混合反对易子为 \(2H\)。既可以使用整体取正的约定，也可以使用其一致共轭的替代约定。是计算，而不是关于符号 \(i\) 的视觉直觉，确定了相对符号。}
\end{classroomqa}

教室里反复核对符号，这种困难成了一个笑话：符号似乎在通勤上班的路上就变了。这个笑话很贴切。让一个格拉斯曼参数穿过一个坐标进行反对易，即使没有任何物理约定改变，也可能改变一个符号。式~\eqref{eq:l7-differential-supercharges}是随后一切内容的固定参照。

\subsection{作为坐标平移的超对称性}

设 \(\xi,\bar\xi\) 为无穷小奇参数。在上述约定下，用
\begin{equation}
\delta_\xi=\xi q+\bar\xi\bar q
\end{equation}
作用在一个超场上，等价于无穷小坐标变换
\begin{align}
\theta&\longrightarrow\theta+\xi,\\
\bar\theta&\longrightarrow\bar\theta+\bar\xi,\\
t&\longrightarrow t+i\left(\xi\bar\theta+\bar\xi\theta\right).
\label{eq:l7-coordinate-shift}
\end{align}
因此标量超场的变化为
\begin{equation}
\delta_\xi\Phi
=\left(\xi\partial_\theta+\bar\xi\partial_{\bar\theta}
+i\xi\bar\theta\partial_t+i\bar\xi\theta\partial_t\right)\Phi
=\left(\xi q+\bar\xi\bar q\right)\Phi.
\end{equation}
这就是超对称性的几何内容。它是奇方向上的平移，并与普通时间中的平移相联系。两个奇平移之所以闭合为一个时间平移，是因为式~\eqref{eq:l7-differential-algebra}。

\begin{classroomqa}
\audiencequestion{有争议的负号是否可能代表时间上的向后运动？\lecturetimestamp{00:58:13}}
\lecturerresponse{不同的符号约定可以反转所显示的无穷小时间平移，但并没有引入一个单独的物理上的逆时过程。重要的是，坐标变换和微分生成元服从同一代数。}
\end{classroomqa}

\begin{classroomqa}
\audiencequestion{这个扩展的时间有物理实现吗，还是只是形式上的记账？\lecturetimestamp{00:58:23}}
\lecturerresponse{在一个有用的意义上，它是记账：它把代数打包成几何。但当引入相互作用时，这种打包就变得具有物理价值。超空间提供了一套系统的演算，用于构造在每一步都保持玻色子--费米子对称性的拉格朗日量，正如张量演算帮助我们构造洛伦兹标量一样。}
\end{classroomqa}

\section{超场与辅助分量}

因为只有两个奇坐标，偶超场的展开会终止：
\begin{equation}
\Phi(t,\theta,\bar\theta)
=\phi(t)+\bar\theta\psi(t)+\bar\psi(t)\theta
+D(t)\bar\theta\theta.
\label{eq:l7-superfield}
\end{equation}
偶单项式的系数 \(\phi\) 和 \(D\) 是玻色性的；奇单项式的系数 \(\psi\) 和 \(\bar\psi\) 是费米性的。没有 \(\theta^2\)、\(\bar\theta^2\) 或更高阶项存活下来。我们让时间依赖保持未展开，因为关于 \(t\) 的普通泰勒级数不会终止，而且在这里也没有用处。

\begin{classroomqa}
\audiencequestion{最高项中的分量 \(D(t)\) 是什么？\lecturetimestamp{00:41:46}}
\lecturerresponse{在这个阶段，它只是有限展开中最后的玻色性系数。在分量作用量中，它通常没有动能项，并通过代数运动方程被消去，因此被称为辅助场。本讲开玩笑地建议用 ``D 表示 dummy（哑变量），''这抓住了它的非动力学角色，但并不是历史上的定义。}
\end{classroomqa}

超场是一个对象，而 \(\phi,\psi,\bar\psi,D\) 是它的分量。超对称变换将它们混合起来。因此，一个玻色子和一个费米子相关联，其结构意义与矢量的 \(x\) 分量和 \(y\) 分量通过旋转相关联大致相同，尽管格拉斯曼方向不能被想象成普通的坐标轴。

\section{超作用量与最高分量}

为了得到运动方程，我们需要一个作用量。在普通场论中，人们构造一个标量拉格朗日量，并把它对时空积分。这里，类似的表达式是
\begin{equation}
S=\int dt\,d\theta\,d\bar\theta\;\Lambda
\bigl(\Phi,\text{ superspace derivatives of }\Phi\bigr),
\label{eq:l7-superaction}
\end{equation}
其中 \(\Lambda\) 是一个偶超场。Berezin 积分是代数性的。我们选择
\begin{equation}
\int d\theta\,d\bar\theta\;\bar\theta\theta=1.
\end{equation}
因此，如果
\begin{equation}
\Lambda=A+\bar\theta B+\bar B\theta+C\bar\theta\theta,
\end{equation}
那么
\begin{equation}
\int d\theta\,d\bar\theta\;\Lambda=C.
\label{eq:l7-top-component}
\end{equation}
奇积分提取出顶分量，留下一个普通拉格朗日量，以便对时间积分。

为什么结果是不变的？超对称变分平移积分变量。顶分量的变分是时间中的普通全导数，而奇全导数在 Berezin 积分下消失。在适当的端点条件下，作用量不变。这是先构造洛伦兹标量、再对时空积分这一做法的超空间类比。

\section{超场的乘法}

从一个旧超场构造新超场的最简单方法是把它与自身相乘。从式~\eqref{eq:l7-superfield}出发，每一项含有 \(\theta^2\) 或 \(\bar\theta^2\) 的项都消失。保持顺序 \(\bar\theta\theta\)，并把带横杠的分量场写在左边，得到
\begin{align}
\Phi^2
={}&\phi^2
+2\phi\bar\theta\psi
+2\phi\bar\psi\theta\notag\\
&+\left(2\phi D-2\bar\psi\psi\right)\bar\theta\theta.
\label{eq:l7-superfield-square}
\end{align}

\begin{figure}[H]
\centering
\includegraphics[width=0.82\textwidth]{lecture_07_figure_05.png}
\lectureframe{01:07:57.844}
\caption{逐分量的超场乘积及其 \(\bar\theta\theta\) 项集合。}
\end{figure}

两项 \(\phi D\) 和 \(D\phi\) 相加，因为两个场都是偶的。两个费米子交叉项在各自重新排序后也相加，但它们的共同符号是负的：
\begin{equation}
(\bar\theta\psi)(\bar\psi\theta)
+(\bar\psi\theta)(\bar\theta\psi)
=-2\bar\psi\psi\,\bar\theta\theta.
\end{equation}
把平方对奇坐标积分得到
\begin{equation}
\int d\theta\,d\bar\theta\;\Phi^2
=2\phi D-2\bar\psi\psi.
\end{equation}
这一项按构造是超对称的，但它还不是动能拉格朗日量。它不包含任何时间导数。

接着，黑板提出 \(\Phi^3\)、\(\Phi^4\) 和 \(\Phi^n\) 作为练习。有一个紧凑的规则。对任意普通函数 \(W\)，
\begin{equation}
\left.W(\Phi)\right|_{\bar\theta\theta}
=W'(\phi)D-W''(\phi)\bar\psi\psi.
\label{eq:l7-function-top-component}
\end{equation}
对于 \(W(\phi)=\phi^2\)，这重现式~\eqref{eq:l7-superfield-square}；对于 \(W(\phi)=\phi^n\)，它给出
\begin{equation}
\left.\Phi^n\right|_{\bar\theta\theta}
=n\phi^{n-1}D-n(n-1)\phi^{n-2}\bar\psi\psi.
\end{equation}
\editorialnote{讲座把 \(n\) 次幂留作练习，并说明预期会出现哪些类型的项。式~\eqref{eq:l7-function-top-component}在固定的排序约定下给出了完整的规则。}

此时，讲师问听众是否已经疲惫。听众说没有；他回答说他自己累了。这段交流不只是插科打诨。超场的幂很容易，因为展开会终止，但导数和约束才是符号记账真正困难的地方。

\section{哪些导数尊重超对称？}

在跟随下一个黑板论证之前，需要作一处更正。普通时间导数是无害的：
\begin{equation}
[\partial_t,q]=[\partial_t,\bar q]=0.
\end{equation}
因此，只要 \(\Phi\) 是超场，\(\partial_t\Phi\) 就按超场变换。微妙的对象是裸奇导数 \(\partial_\theta\Phi\) 和 \(\partial_{\bar\theta}\Phi\)，因为它们不与两个超荷都反对易。我们需要修改后的奇导数，用于约束以及明显超对称的分量构造。
\editorialnote{在这里使用的约定中，关于时间微分会无法产生另一个超场的口头说法并不正确。由于 \(\partial_t\) 与两个超荷都对易，它是协变的。随后构造修改后的奇导数仍然是必要的，而这正是本讲座真正的数学要点。}

\subsection{``魔场''类比}

讲座用普通旋转来启发这个判据。设 \(f(x,y)\) 满足
\begin{equation}
\partial_y f=0.
\label{eq:l7-magic-constraint}
\end{equation}
称这样的函数为 ``魔场。'' 它依赖于 \(x\)，但不依赖于 \(y\)。魔场这一类在旋转下是否保持？对于
\begin{equation}
L_z=-i\left(x\partial_y-y\partial_x\right),
\end{equation}
我们发现
\begin{equation}
[\partial_y,L_z]=i\partial_x\neq0.
\end{equation}
因此，即使 \(f\) 最初满足式~\eqref{eq:l7-magic-constraint}，\(L_z f\) 一般也会获得 \(y\) 依赖性。该约束不是旋转不变的。

一般教训很简单。如果一个场满足 \(C\Phi=0\)，而一个对称性由 \(G\) 生成，那么当 \(C\) 与 \(G\) 的分次对易子消失时，变换后的场仍然受约束。对于两个奇算符，分次对易子就是反对易子。

\begin{classroomqa}
\audiencequestion{一个关于场的约束何时与一个对称性相容？\lecturetimestamp{01:17:56}}
\lecturerresponse{当约束算符与对称性生成元分次对易时。于是 \(C(G\Phi)\) 可以被移过 \(G\)，并约化为 \(G(C\Phi)=0\)。旋转的例子失败，是因为 \(\partial_y\) 不与旋转生成元对易。}
\end{classroomqa}

\subsection{协变奇导数}

在超荷由式~\eqref{eq:l7-differential-supercharges}固定后，定义
\begin{equation}
\mathcal D=\partial_\theta-i\bar\theta\partial_t,
\qquad
\bar{\mathcal D}=\partial_{\bar\theta}-i\theta\partial_t.
\label{eq:l7-covariant-derivatives}
\end{equation}
唯一的变化是时间导数项的符号。直接计算给出
\begin{align}
\{\mathcal D,q\}
&=\{\mathcal D,\bar q\}
=\{\bar{\mathcal D},q\}
=\{\bar{\mathcal D},\bar q\}=0,\\
\{\mathcal D,\bar{\mathcal D}\}&=-2i\partial_t.
\label{eq:l7-covariant-algebra}
\end{align}
两个可能非零的交叉项相互抵消。这是协变导数的奇类比：它在微分的同时保持变换规律。

讲义中\(D\)既用于辅助分量，也用于这个导数。我们保留\(D(t)\)表示分量，并用花体\(\mathcal D\)表示算符。

\section{手征约束}

现在我们可以施加一个在超对称下仍然成立的非平凡约束：
\begin{equation}
\bar{\mathcal D}\Phi=0.
\label{eq:l7-chiral-constraint}
\end{equation}
这样的场称为手征场。如果它在一次超对称变换之前是手征的，那么
\begin{equation}
\bar{\mathcal D}(q\Phi)
=-q(\bar{\mathcal D}\Phi)=0,
\qquad
\bar{\mathcal D}(\bar q\Phi)
=-\bar q(\bar{\mathcal D}\Phi)=0.
\end{equation}
因此整个变换后的超场仍留在受约束的类中。与魔场条件不同，手征性与该对称性相容。

为求解该约束，定义一个平移后的时间坐标
\begin{equation}
\tau=t+i\bar\theta\theta.
\label{eq:l7-chiral-time}
\end{equation}
由于我们使用左导数，
\begin{equation}
\partial_{\bar\theta}\tau=i\theta,
\end{equation}
因此
\begin{equation}
\bar{\mathcal D}F(\tau,\theta)
=\left(\partial_{\bar\theta}-i\theta\partial_t\right)F(\tau,\theta)
=i\theta\partial_\tau F-i\theta\partial_\tau F=0.
\end{equation}
所以一般的局域手征超场仅通过\(\tau\)依赖于\(\bar\theta\)：
\begin{equation}
\Phi_{\mathrm{ch}}(\tau,\theta)
=\phi(\tau)+\bar\psi(\tau)\theta.
\label{eq:l7-chiral-field-tau}
\end{equation}

\begin{classroomqa}
\audiencequestion{仅依赖于\(\tau\)和\(\theta\)是满足手征约束的唯一形式吗？\lecturetimestamp{01:32:35}}
\lecturerresponse{局域地看，是的。微分方程\(\bar{\mathcal D}\Phi=0\)恰恰说明，独立的\(\bar\theta\)依赖可以被吸收进平移后的坐标\(\tau=t+i\bar\theta\theta\)。剩下的有限展开是\(\tau\)的函数，再加上一个正比于\(\theta\)的项。}
\end{classroomqa}

用原来的时间坐标重新展开，
\begin{align}
\phi(\tau)
&=\phi(t)+i\bar\theta\theta\dot\phi(t),\\
\bar\psi(\tau)\theta
&=\bar\psi(t)\theta,
\end{align}
因为第二行的修正包含\(\theta^2\)。因此
\begin{equation}
\Phi_{\mathrm{ch}}(t,\theta,\bar\theta)
=\phi(t)+\bar\psi(t)\theta
+i\bar\theta\theta\dot\phi(t).
\label{eq:l7-chiral-field-components}
\end{equation}

\begin{classroomqa}
\audiencequestion{\(\phi(t+i\bar\theta\theta)\)的泰勒展开是否应包含二阶项？\lecturetimestamp{01:33:51}}
\lecturerresponse{不。第一项修正是\(i\bar\theta\theta\dot\phi\)，而每一项更高阶修正都包含\((\bar\theta\theta)^2=0\)。听众立刻抓住了这一点；幂零性已经提供了前两项并终止了级数。}
\end{classroomqa}

这就是手征性的经济性。式~\eqref{eq:l7-superfield}中未受约束的超场有两个玻色分量槽和两个费米分量槽。手征条件在这个最小的一维例子中去掉了独立的辅助槽，并将最高分量与\(\dot\phi\)联系起来。

\section{恰当停止的计算}

自然的下一步是将手征场与其共轭结合起来，并在超空间上积分。共轭展开开始了，一位听众问，在共轭下\(\theta\bar\theta\)的次序是否必须反转。此时符号已不再可信。与其虚张声势，计算就此停止。

\begin{classroomqa}
\audiencequestion{即使推迟共轭场的乘法，我们能否先得到关键结论？\lecturetimestamp{01:37:30}}
\lecturerresponse{可以。在符号和共轭被一致处理后，由手征场及其共轭构成的超空间表达式会产生所需的普通玻色和费米动能项。本讲的重要结果是构造原理，而不是从疲惫计算中抄来的不可靠的最后一行。}
\end{classroomqa}

停止点应当保持可见。这套形式体系毫不宽容：一个放错位置的奇因子就会改变一个符号，而同一套约定必须在共轭、微分、乘法和积分中始终成立。失去这条链条时，正确的回应是先修复它，再继续。在标准约定中，熟悉的动能作用量由共轭乘积构成，但其详细的分量约化属于下一次干净的计算，而不是对未完成的板书行进行重建。

\section{为什么对称性才是关键结论}

超对称拉格朗日量所做的不仅仅是把一个玻色子和一个费米子并排放置。它约束它们的质量和相互作用，使它们构成超代数的一个表示。在严格超对称的真空中，同一超多重态的成员具有相等的质量。它们的圈贡献相互关联，玻色和费米的紫外效应在对称性所强制的精确组合中相消。

这就是构造机制之所以重要的原因。人们可以列出费曼顶点并逐图验证相消，但那样相消看起来就像是偶然的。超对称解释它们的方式，正如洛伦兹对称性解释为什么散射振幅在每个惯性系中都一致地变换。

\editorialnote{精确的超对称性并不意味着每一个单独的自能修正都会消失。它关联的是重整化，并消除特定的紫外敏感性；非重整化陈述取决于所讨论的量、维数以及超对称的数量。本讲中关于广泛抵消的说法以这种限定形式保留。}

\begin{classroomqa}
\audiencequestion{``chiral''怎么拼写？\lecturetimestamp{01:39:48}}
\lecturerresponse{C-H-I-R-A-L。这个词标记的是约束\(\bar{\mathcal D}\Phi=0\)；它并不是三维空间中物体的普通几何手性，尽管这一术语源自相对论性旋量的左手和右手结构。}
\end{classroomqa}

\begin{classroomqa}
\audiencequestion{能否绕过拉格朗日量，直接得到超对称费曼规则？\lecturetimestamp{01:40:00}}
\lecturerresponse{当然可以陈述这些规则，但那样一来，耦合之间的关系以及它们所产生的抵消就会显得像是无法解释的巧合。更深层的内容是广义时空对称性。拉格朗日量或等价的不变形式使这一原因变得可见，并系统地生成这些规则。}
\end{classroomqa}

假设每一个微扰阶都碰巧给出洛伦兹不变的振幅，尽管从未提到洛伦兹变换。经过足够多的计算之后，我们会要求一个解释。答案将是那个一次性组织所有阶的时空对称性。超空间的作用与此类似：\(\theta,\bar\theta\)中的平移混合玻色子和费米子分量，而它们的代数在时空平移上闭合。

\begin{classroomqa}
\audiencequestion{普通空间坐标是否也应像格拉斯曼坐标那样展开成有限幂级数？\lecturetimestamp{01:44:55}}
\lecturerresponse{不。普通坐标的幂级数不会终止，通常也不是自然的分解。傅里叶变换更有用，因为它把场分解为波和动量。按\(\theta\)展开之所以特殊，正是因为幂零性使其成为有限展开。}
\end{classroomqa}

\begin{classroomqa}
\audiencequestion{如果LHC出现了一个新粒子，我们如何判断它是超伴侣，而不只是一个未知粒子？\lecturetimestamp{01:45:42}}
\lecturerresponse{超对称预言的是相互关联的量子数和相互作用，而不只是一个额外态。自旋相差二分之一，规范荷相互关联，许多三次和四次耦合由比一般理论更少的独立参数决定。截面、分支比以及谱的其余部分都必须满足这些关系。}
\end{classroomqa}

一个不熟悉的共振本身并不能证明超对称。一个由耦合关系和谱关系构成的网络才能证明。这就是超空间构造在实验上的对应物：表面上独立的分量相互作用源自同一个不变对象。

\begin{classroomqa}
\audiencequestion{观测到的粒子及其假想超伴侣并不具有相等的质量。这会如何改变理论？\lecturetimestamp{01:47:01}}
\lecturerresponse{改变极大。超对称必须破缺，通常通过自发破缺并传递到观测部门来实现。破缺使多重态质量分裂，同时保留足够多的结构以控制相互作用，并且如果标度足够低，也控制紫外修正。弄清这一机制本身就是一套庞大的理论，尤其是在包含引力时。}
\end{classroomqa}

\section{原理简单，计算困难}

结尾的评价异常坦率。出发点听起来很简单：玻色子和费米子由一种对称性联系起来。然而具体实现很快就变成一门关于指标、共轭、格拉斯曼排序和符号的学问。即使在构造出这种对称性之后，现实的粒子物理学仍需要破缺它；把它与引力耦合会引向超引力，并增加另一层结构。

这种困难并不证明这个想法是空洞的。大自然没有义务让一个简单的原理易于计算。专家们通过反复处理细节、直到符号规则成为第二天性来学习这个课题。挫折感来自紧凑的代数
\begin{equation}
\{Q,Q^\dagger\}=2H
\end{equation}
与在相互作用理论中实现它所需的大量机械之间的差距。

\begin{classroomqa}
\audiencequestion{另一层抽象最终会简化形式体系，还是复杂性会不断累积？\lecturetimestamp{01:49:33}}
\lecturerresponse{还没有任何一般性的简化能免除这些细致工作。底层规则可能很短，但要可靠地使用它们，就需要反复掌控约定。这就是为什么完整的处理通常要占整整一门技术课程，而不是宽泛粒子物理学课程的一小部分。}
\end{classroomqa}

\begin{classroomqa}
\audiencequestion{这些变换是否告诉我们费米子和玻色子究竟是什么，或者为什么它们的日常行为不同？\lecturetimestamp{01:51:42}}
\lecturerresponse{它们告诉我们，二者是同一个超场的分量，并且在超荷作用下互为映像。用\(Q\)作用会把一个玻色型分量映射为一个费米型分量，而下一次变换则产生导数和辅助分量。这类似于旋转混合矢量分量，但不能被翻译成一幅同样字面的图像：恰当的语言正是格拉斯曼几何本身。}
\end{classroomqa}

这个回答把我们带回核心思想。超对称变换是超空间中的微小运动。它混合\(\phi,\psi,\bar\psi,D\)，正如普通旋转混合分量一样，并且在这样做的同时闭合到时间平移上。这幅图像是抽象的，但并非任意的：超场中的每一项、协变导数中的每一个符号，以及分量相互作用之间的每一种关系，都是为了实现该代数而存在的。

\section{我们构建了什么}

我们从匹配的振子哈密顿量出发，并恢复了它的两种不同的拉格朗日描述。玻色子服从二阶振子方程；费米子则需要一阶格拉斯曼作用量以及一对复的奇分量。超荷交换这两个模式，各自平方为零，反对易得到哈密顿量，并与它交换。

把这些荷表示为微分算子引出了超空间。有限的超场展开汇集了玻色型、费米型和辅助分量；别列津积分选出超拉格朗日量的顶分量；而乘法生成了不变的相互作用候选者。随后，对与对称性相容的奇微分的要求引出了\(\mathcal D,\bar{\mathcal D}\)、手征约束以及平移后的坐标\(\tau=t+i\bar\theta\theta\)。

未完成的共轭乘积计算是这一课的一部分，而不是需要隐藏的缺陷。超对称始于一个简单的物理要求，但它的微积分对每一种排序选择都很敏感。一旦这些选择得到控制，回报就是强大的：质量和耦合由对称性组织起来，量子修正满足并非偶然的关系，而玻色子和费米子成为一个几何对象的分量。
